% AI-produced by GPT-6 Astra, OpenAI, 20 November 2026.
% Authored scope note only: not numbered, promoted, or included in the paper.
\subsection*{The ordinary rank-four formal comparison has no extra fourth homotopy group}

Use the standard classifying-space homotopy-fibre conventions
\[
 X=\operatorname{hofib}(BO(4)\longrightarrow BPL(4)),\qquad
 Y=\operatorname{hofib}(BO(4)\longrightarrow BTOP(4)),\qquad
 Z=\operatorname{hofib}(BPL(4)\longrightarrow BTOP(4)).
\]
These are the ordinary, non-equivariant spaces customarily written
$PL_4/O_4$, $TOP_4/O_4$, and $TOP_4/PL_4$. Their maps are the standard
maps of classifying spaces; no literal inclusion of arbitrary singular
orthogonal families in the jointly PL simplicial group is asserted.
The triangle of classifying-space maps gives the based fibre sequence
\[
 X\longrightarrow Y\longrightarrow Z.
\]

The inspected Burghelea--Lashof II original, Section5, printed p.37,
formula(3), gives $\pi_iX\cong\pi_i(PL/O)$ for $i<4$; its precise use
and finite-rank convention are recorded in A.128, equations(A.128.1--2),
and A.132. The stable groups vanish for $0\leq i\leq3$. Thus $X$ is
3-connected. A.132 uses the corrected finite-rank TOP theorem to show
that $Y\to TOP/O$ is 5-connected. Kirby--Siebenmann, EssayV, Theorem5.5,
printed p.251, then gives
\[
 \pi_iY=0\quad(i=1,2,4),\qquad \pi_3Y\cong\mathbb F_2.
\]
The fibre sequence therefore gives $\pi_iZ=0$ for $i=1,2$ and
$\pi_3Z\cong\mathbb F_2$. Crucially, its degree-four segment is
\[
 \pi_4Y=0\longrightarrow\pi_4Z
       \longrightarrow\pi_3X=0.
\]
Consequently $\pi_4Z=0$; no knowledge of $\pi_4X$ is needed. The spaces
are connected in these ordinary comparison conventions: the connectedness
of $Y$ makes the orthogonal component map onto the TOP component group,
and that map factors through the PL component group. Thus the ordinary
4-type of $Z$ is $K(\mathbb F_2,3)$. This argument is only a finite-rank
homotopy calculation from the already inspected inputs.

\textbf{Exact scope.}\quad
This does not identify an equivariant quotient for the sign representation
$4\sigma$ of $C_2$. In particular it computes neither the odd joint-germ
forgetting kernel in A.162 nor the PL image in that kernel. It also does
not turn a formal rank-four reduction on a four-dimensional compact
cylinder into an isotopy of actual structures. The missing statement in
the RP3 zero-Casson--Sullivan question remains a relative, proper geometric
realization statement. The KS open-four sliced classification and its
relative end conditions cannot be replaced by this homotopy calculation;
the required compact collar/orbit step remains unproved. There is no new
numbered result, no original-row closure, and no qualifying gate movement.
